\begin{center}
	\Large{\textbf{EXECUTIVE SUMMARY}}	
\end{center}

\vspace{1em}

\normalsize

Student academic burnout is an increasingly critical concern in higher education, characterized by emotional exhaustion, reduced academic performance, and disengagement from learning activities. Early detection is essential for timely intervention; however, most existing approaches rely on subjective self-reporting rather than objective behavioral signals.

This project presents an early warning machine learning system designed to detect academic burnout risk using multi-source student behavioral data. The system integrates data from multiple categories including academic performance, attendance, assignment submission patterns, Learning Management System (LMS) engagement, and wellbeing indicators to construct a comprehensive student profile with engineered features.

Multiple machine learning models such as Random Forest, Logistic Regression, and Decision Tree were trained and evaluated for burnout prediction. Among these, Random Forest achieved the best performance with high accuracy, precision, and recall. Key indicators contributing to burnout prediction include declining academic performance, reduced engagement, irregular attendance, and increased stress-related patterns.

The system assigns each student a risk score, which is further categorized into low, medium, and high risk levels to support early intervention. The results demonstrate that behavioral data can effectively identify burnout patterns, providing a scalable and practical solution for improving student well-being and academic outcomes.